\scanpage{82}
\margin{81}
\texttt{D A T A S E T S\quad P I C T U R E S}
\begin{verbatim}
******************************************************************************************************************
*              *                  *        *       *             *            *     *       *        *           *
*   DDNAME     * DESTINATION      * CLASS  *COPIES *  RECORDS-NB *  BYTES-NB  *FORM * FCB   * UCS    *SPIN HOLD  *
*              *                  *        *       *             *            *     *       *        *           *
******************************************************************************************************************
*              *                  *        *       *             *            *     *       *        *           *
*    SYSOS3    * CNU     M06      *    A   *   1   *         67  *      3,893 *STD. * ****  * ****   *           *
******************************************************************************************************************
\end{verbatim}

\deleted{\diagram{\begin{array}{ccc@{\qquad}ccc@{\qquad}ccc}
\chi_0&\chi_1&\chi_2&&&&&&\\
1&\bar\epsilon&0&1&1&0&1&1&0\\
1&\unsure{\xi}&0&1&\unsure{\xi}&0&1&\unsure{0|\xi}&\unsure{0}\\
1&\unsure{\eta}&1&1&\unsure{\eta}&1&1&\unsure{\eta}&\unsure{1}\\
&\unsure{\xi}-1&&&&&1-\unsure{\bar\epsilon}=1&&
\end{array}}}

\deleted{\diagram{\begin{array}{ll}
\left.\begin{array}{ll}\chi_0&\text{orientation de }D\\
\chi_1&\text{sens des }(n\text{-gones}\ \inserted{\text{duaux}}\text{ inscrits})\end{array}\right\}
&\begin{array}{l}\text{caractères fondamentaux}\\\text{sur }\unsure{\mathbb D_2}\end{array}\\
\chi_2&\text{caractère parité}
\end{array}}}

\deleted{\uline{NB} le caractère d'orientation pour
$\unsure{\mathcal X}$ est $\chi_0\chi_2$,
le caractère ``\unsure{longueurs circuits}''
est $\chi_0\chi_1$}

\deleted{\[\epsilon=\text{parité de }n\in\{0,1\}\qquad
\bar\epsilon=\text{opposé de }\epsilon=\epsilon-1\bmod2\]}

\deleted{\diagram{\begin{array}{lccc}
&\chi_0&\chi_1&\chi_2\\
\begin{array}{l}\frac12\text{ tour }\deleted{\unsure{autour}}\text{ d'un trou}\\
(\text{en }2n+\uuline{2}\text{ pas})\end{array}
&1&\unsure{\bar\epsilon}\deleted{\unsure{1}}&0\\
\begin{array}{l}\frac12\text{ tour autour}\\\deleted{\unsure{d'un sommet}}\\\deleted{\unsure{d'un point}}\\\text{d'ordre }\alpha\end{array}
&1&?&\alpha
\end{array}}}

\deleted{$\ell$ lacet en $F$}

\deleted{\deleted{$\unsure{\beta}$} $\unsure{u}(\ell)$ : Autom de $\operatorname{Pol}_0(F)$ (polygone de degré $2n$)}
\deleted{\[\chi_0(\ell)\mathrel{\mathop=^{\mathrm{def}}}\unsure{\operatorname{sgn}}\,\unsure{u}(\ell)=\left\{\vphantom{\frac11}\right.\]}
\deleted{\deleted{\unsure{si de}}}

\deleted{\deleted{1) Si $\alpha$ pair alors $\chi_0$ et}}

\deleted{Prenant $\alpha$ \deleted{pair} \inserted{impair}, on voit que
\uline{$\chi_0,\chi_2$ indépendants} / et $\chi_1,\chi_2$ aussi si $n$
pair, i.e. $\bar\epsilon=1$
\[\text{i.e. }\chi_0,\unsure{\chi_2|\chi}\text{ indépendants}\]}

\deleted{\uline{$\chi_0$}\qquad \uline{$n$ pair}\qquad
$\Pi\longrightarrow\mathbb D_{2n}$ surjectif
\[\text{et : }\Pi\longrightarrow\mathbb D_{2n}\times\{\pm1\}\text{ a une image qui contient }\mathbb D_n\times\{\pm1\}\]}

\deleted{\diagram{\begin{array}{lccc}
\uline{n\text{ pair}}&1&1&0\\
&1&?&\alpha\\[1em]
\uline{n\text{ impair}}&1&0&0\\
&1&?&\alpha
\end{array}}}

\deleted{$\chi_2$ n'est pas lié à $\chi_0$, ni à $\chi_2$ --- il est possible que
$\chi_2,\chi_1$ le soient, d'où tjrs $?=1$}
\deleted{\margin{$\chi_0$}}
\figure{\begin{tikzpicture}[x=.006cm,y=-.006cm]
\draw (0,100)--(1000,1800);
\draw (600,0)--(1800,1700);
\draw (-500,700)--(1700,1900);
\draw (650,1850)--(650,2050);
\draw (650,1850)--(1700,1820) .. controls (1830,1810) and (1790,1740) .. (1840,1730)--(2300,1730);
\end{tikzpicture}}
